\section{Limitations}\label{sec:limits}
\paragraph{Reference and numerical scope.} Linear optimality is conditional on the specified covariance, centring, operator and noise. The image-periodogram oracle is truth-dependent. E3's assumed variance, clipping, amplitude scaling and missing pixel/energy/residual archives limit its use to the recovered descriptive comparison. Exact matrix identities do not certify an approximate solver or a clipped output. E1 also has a native forward/inverse discretisation mismatch.

\paragraph{Empirical generalisation.} ASISR results cover a limited set of decoders, training recipes, datasets and three seeds. Bands and cells share images and are dependent. Similar recovery is not a Bayes limit or proof that more capable decoders cannot improve. Nyquist normalisation improves consistency but fails the original strict threshold in Set14; content sensitivity and readout losses remain empirically observed, not causal explanations.

\paragraph{Mechanism and guidance.} Classical scalar identities require fixed effective noise and cannot be substituted into a nonlinear network without further assumptions. G1 supports only coarse ordering; L2 fails function similarity and general superiority over a simpler predictor. Threshold sensitivity at $.1/.8$ remains untested. G2 covers one operator, one calibration and fixed noise draws; its constant selection and worse comparison to CV leave adaptive and computational benefits unsupported. No learned INR model-choice rule is established.

\paragraph{Preserved auxiliary evidence.} Supplementary controls use already-seen development images and single seeds. Their reverse rankings and static-baseline equivalence are retained. The Gaussian ambiguity example depends on the allowed estimator class; a restricted two-candidate regret bound is not an all-estimator reconstruction lower bound. Classical covariance-estimation and design tools are not claimed as new theorems. Calibration-only EEG does not supply independent cross-modal validation.

\section{Conclusion}
The paper baseline concerns scale-dependent frequency recovery and readout effects in ASISR. Nyquist normalisation reduces scale variation, and coarse readout-loss ordering replicates across two architecture pairs and two additional seeds. Stronger collapse, function-equivalence, quantitative mechanism and adaptive-selection claims remain unsupported by their respective tests. Covariance-conditioned linear references and classical mismatch costs help organise these observations, with their assumptions and numerical limits stated separately. Further innovation should be assessed against this preserved baseline rather than inferred from additional correlations or oracle comparisons.
